\begin{claim}{D.70}{Proved}{The quartic torus multiplier is globally one}
Every element of order four in the full loop-outer group
of the once-punctured torus has multiplier
$1\in\widehat\Z^\times$.
Consequently every finite two-subgroup of that full group
has multiplier image in the two global units $\{1,-1\}$,
and the multiplier is determined by its $2$-adic coordinate.
The conclusion concerns the full group, without a
geometric-membership assumption.
It determines the quartic multiplier; the conjugacy of
an arbitrary quartic action to a specified geometric
action is not asserted.
\end{claim}
\begin{proof}
\emph{An exact degree-three restriction obstruction.}
We first give the cohomological calculation used below.
Let $P$ be the pro-two completion of an orientable closed
surface group of genus $g\geq2$, and let $b$ be an actual
involution of $P$ whose action on $H^1(P,\F_2)$ is
identity and whose action on $H^1(P,\Z_2)$ is nonidentity.
Write
\[
 G=P\rtimes_b\langle\tau\rangle,\qquad
 \tau^2=1,\qquad
 u:G\longrightarrow\F_2,\quad u(h\tau^\epsilon)=\epsilon .
                                                               \tag{D.70a}
\]
Then any two $b$-fixed continuous characters
$\theta,\psi:P\to\Z_2$ have zero reduced cup product
in $H^2(P,\F_2)$.

Here are all maps and hypotheses for this assertion.
The group $P$ is a torsion-free Demushkin pro-two group:
its cohomological dimension is two, its first cohomology
has dimension $2g$, and its second cohomology is
$\F_2$ with the nondegenerate surface cup pairing.
The integral character module is $\Z_2^{2g}$.
These are the surface duality comparisons of A.2 and A.7,
used with both completions in D.42.

Every $\xi\in H^1(P,\F_2)$ has the actual character
extension
\[
 \widetilde\xi:G\longrightarrow\F_2,\qquad
                  \widetilde\xi(h\tau^\epsilon)=\xi(h),
                                                               \tag{D.70b}
\]
because $b$ fixes $\xi$.
The action on $H^2(P,\F_2)$ is also identity, since this
space is one-dimensional over $\F_2$.
The multiplicative Hochschild--Serre spectral sequence
therefore has its actual second page
\[
 E_2^{i,j}
   =H^i(C_2,\F_2)\otimes H^j(P,\F_2)
   =\F_2[u]\otimes H^j(P,\F_2),\qquad 0\leq j\leq2 .
                                                               \tag{D.70c}
\]
The quotient class $u$ is permanent: the section
$C_2\to G$, $\tau^\epsilon\mapsto\tau^\epsilon$,
makes inflation injective.
All classes in $E_2^{0,1}$ are permanent by (D.70b).
The surface class in $E_2^{0,2}$ is the cup product of
two of them, so it is permanent as well.
These classes generate the entire second-page algebra.
Thus every differential vanishes, and multiplication by
$u$ on its row-two term gives the nonzero class
$u\eta\in E_\infty^{1,2}$ for each nonzero surface
class $\eta$.

The fixed integral characters have the stronger actual
extensions
\[
 \widetilde\theta:G\longrightarrow\Z_2,\qquad
 \widetilde\psi:G\longrightarrow\Z_2,\qquad
 \widetilde\theta(h\tau^\epsilon)=\theta(h),\quad
 \widetilde\psi(h\tau^\epsilon)=\psi(h).
                                                               \tag{D.70d}
\]
Every finite subgroup of $G$ injects into $G/P=C_2$,
because $P$ is torsion-free. For each order-two
subgroup $T=\langle z\rangle$, the two characters
in (D.70d) vanish on $T$: twice their values on
$z$ is zero in the torsion-free group $\Z_2$.
There is the exact direct-product normalizer
\[
 \begin{aligned}
 C_P(T)\times T&\longrightarrow N_G(T),
                    &(h,z^\epsilon)&\longmapsto hz^\epsilon,\\
 N_G(T)&\longrightarrow C_P(T)\times T,
                    &w&\longmapsto(wz^{-\epsilon},z^\epsilon),
                    \qquad \epsilon=u(w).
 \end{aligned}                                                \tag{D.70e}
\]
Indeed every normalizer element fixes the unique
nonidentity element of $T$, and its quotient label
determines $\epsilon$ uniquely.
These are inverse continuous homomorphisms.

Write $z=n\tau$. Its actual action on $P$ is
$\operatorname{Inn}(n)b$, with square identity.
It has the same integral homology action as $b$,
so it is nonidentity by the stated hypothesis.
The hypotheses of
\cite[Corollary~2 and its proof in \S2]{A-KochloukovaZalesskii2007}
therefore hold for this particular action:
$P$ is a pro-two Poincare duality group of dimension
two and type $\mathrm{FP}_\infty$, and the acting
subgroup has order two in $\operatorname{Aut}(P)$.
The corollary proves that the full fixed subgroup
$C_P(T)$ is a free pro-two group.
In particular
\[
                          H^2(C_P(T),\F_2)=0.                \tag{D.70f}
\]
For clarity, the infinite-index step in that proof
uses its Lemma~1 for an actual nontrivial finite
two-group of automorphisms of a torsion-free pro-two
group; it then applies the dimension bound for an
infinite-index subgroup of a Demushkin group.
It is applicable separately to every section $T$.

Let bars denote the reductions of (D.70d) modulo two.
Under (D.70e), both restricted characters are pulled
back from $C_P(T)$, since they vanish on $T$.
Their cup product is zero by (D.70f).
It follows that
\[
 \operatorname{res}_{N_G(T)}
           \bigl(u\smile\overline{\widetilde\theta}
                    \smile\overline{\widetilde\psi}\bigr)=0
                 \quad\text{for every }T.                  \tag{D.70g}
\]
The full Brown restriction theorem
\cite[Theorem~2]{A-KochloukovaZalesskii2007}
is injective in degree three here: $P$ is open
of cohomological dimension two, and $G$ has no
$C_2\times C_2$ because its finite subgroups
inject into $C_2$.
Thus the class in (D.70g) is zero in $H^3(G,\F_2)$.
If $\bar\theta\smile\bar\psi=\eta\ne0$ on $P$,
its associated graded class under (D.70c) is
the nonzero element $u\eta\in E_\infty^{1,2}$,
a contradiction. This proves the asserted
zero-cup property of the integral fixed characters.
Only the injectivity part of Brown restriction
is needed in this calculation.

\emph{The actual full cusp filling and its pro-two quotient.}
Now take an order-four full torus action and the
actual lift $\alpha$ of D.46.
D.43 and D.50 give
$\lambda^2=1$ and $\lambda_2=1$.
Suppose an odd coordinate satisfies $\lambda_q=-1$.
D.68 gives its actual dihedral quotient, and D.69
gives the actual type-$(q,2)$ cover
\[
 N=\ker\bigl(\rho:F\twoheadrightarrow D_{2q}\bigr),
 \qquad \beta=\alpha^2,\qquad
 R=\langle\rho(y)\rangle\cong C_q .
                                                               \tag{D.70h}
\]
In particular $q\equiv1\pmod4$, so $q\geq5$.
Let $\Gamma_N$ be the corresponding discrete
finite-index subgroup of the original punctured-torus
group, with its original covering surface.
Use the two actual based peripheral generators
$d_0=c^q$ and $d_1=xc^qx^{-1}$ of this cover.
The second base path is the lift of the original
word $x$ to the other peripheral orbit.
Their full procyclic subgroups in $N$ are the
two covering cusps.
Let
\[
 K=\overline{\langle\!\langle d_0,d_1\rangle\!\rangle}_N,
 \qquad
 \Pi=N/K,\qquad P=\Pi^{(2)},\qquad
 p:N\twoheadrightarrow\Pi\twoheadrightarrow P .
                                                               \tag{D.70i}
\]
These are the specified quotient maps.
Filling the two punctures of the discrete cover
gives the closed genus-$q$ surface group
$\Gamma_N/\langle\!\langle d_0,d_1\rangle\!\rangle$.
Its full completion is exactly $\Pi$:
a map from either group to a finite group is
precisely a map from $\Gamma_N$ killing $d_0,d_1$.
The finite-index subgroup $\Gamma_N$ has its full
induced profinite topology in $F$, as in D.69,
so this identifies that universal finite quotient
with the displayed $N/K$.
Restricting the same finite targets to two-groups
gives the displayed maximal pro-two quotient $P$.
Consequently $P$ is the genus-$q$ pro-two closed
surface group used in the preceding calculation.

The preservation of $K$ has explicit conjugators
in $N$. Retain the cusp root $v=f_0s$ of D.68,
so $\alpha(c)=f_0^{-1}c^\lambda f_0$.
For $g\in\{1,x\}$ choose $g'\in\{1,x\}$
with $\rho(g')R=\rho(\alpha(g)f_0^{-1})R$.
Since $\rho(c)=r^{-2}$ generates $R$, there is
an integer $k$ for which
$\rho(\alpha(g)f_0^{-1})=\rho(g'c^k)$.
The actual element
$n=\alpha(g)f_0^{-1}c^{-k}(g')^{-1}$ then belongs to $N$,
and direct multiplication gives
\[
 \alpha(gc^qg^{-1})
       =n\bigl(g'c^q(g')^{-1}\bigr)^\lambda n^{-1}.
\]
The inverse automorphism gives the reverse inclusion.
Thus $\alpha(K)=K$ with the original two labels
and their full unit powers retained.
The maximal pro-two kernel of $\Pi$ is
characteristic, so $\alpha$, and hence $\beta$,
induce actual automorphisms through both arrows
in (D.70i). Denote the latter induced automorphism
of $P$ by $b$; it satisfies $b^2=1$.
No section is changed in this passage.

\emph{The dual inclusion and the closed integral deck splitting.}
The actual pullback of characters gives an injection
\[
 p^*:L=H^1(P,\Z_2)\hookrightarrow
                     V_2=H^1(N,\Z_2).
                                                               \tag{D.70j}
\]
Its image is exactly the characters which vanish
on $d_0,d_1$.
One direction follows from (D.70i).
Conversely such a character kills $K$, factors
through $\Pi$, and then through $P$ because its
target $\Z_2$ is pro-two.
This proves both directions of the assertion
for the exact displayed maps.
The same argument holds with target $\F_2$,
and reduction identifies $L/2L$ with $H^1(P,\F_2)$.

The rotation deck action on the punctured cover
extends over each of its two filled punctures.
It preserves orientation and gives the actual
action of $R$ on $L$ and its surface cup form.
As in D.69 put
\[
 e=\frac1q\sum_{r'\in R}r',\qquad W_2=(1-e)V_2 .
                                                               \tag{D.70k}
\]
Both cusp restriction functionals on $V_2$
are $R$-invariant. Indeed $R$ fixes each of the
two cusp labels, and its oriented deck action
changes a based peripheral generator only by
conjugation in $N$.
Consequently each such functional vanishes on
$W_2$, so (D.70j) identifies $W_2$ with
$(1-e)L$.

On the deck-trivial part, restriction and transfer
identify $eV_2$ with $H^1(F_0,\Z_2)$, where
$F_0=\rho^{-1}(R)$ is the type-$(1,2)$ double cover
of D.68. Its retained basis is
$u=x^2,B=y,C=xyx^{-1}$.
The original boundary is $c=CB^{-1}$,
and its two covering boundary homology classes
are opposite.
Restriction to $N$ sends each peripheral evaluation
to $q$ times the corresponding evaluation on $F_0$.
Since $q$ is a unit in $\Z_2$, the condition in
(D.70j) on this part is exactly
\[
                              \xi(B)=\xi(C).
                                                               \tag{D.70l}
\]
Thus $eL$ has rank two, and the actual direct
sum of integral character modules is
\[
                  L=eL\oplus W_2,\qquad
             \operatorname{rank}(W_2)=2q-2 .
                                                               \tag{D.70m}
\]
D.68 and the exact dual convention in D.69 give
$\beta^*=-J^{\mathsf T}$ on $eV_2$, with $J$
fixing $u$ and interchanging $B,C$.
On (D.70l), $J^{\mathsf T}$ is identity.
D.69 gives $\beta^*=I$ on $W_2$.
Equivariance of (D.70j) therefore gives the
full integral action
\[
                  b^*|_{eL}=-I,\qquad b^*|_{W_2}=I.
                                                               \tag{D.70n}
\]
It follows that $b$ has order exactly two,
its integral homology action is nonidentity,
and its action on $L/2L$ is identity.
All three assertions follow from the actual
rank-two summand in (D.70n); none is inferred
from the order of the original outer class alone.

\emph{The unimodular fixed cup pair and the contradiction.}
Choose the oriented surface generator of
$H^2(P,\Z_2)\cong\Z_2$.
The cup pairing
\[
               \langle\, ,\,\rangle:L\times L\to\Z_2
                                                               \tag{D.70o}
\]
is alternating and unimodular.
Every deck operator $r'$ preserves this generator,
so
$\langle r'\xi,\eta\rangle
 =\langle\xi,(r')^{-1}\eta\rangle$.
Summing over $R$ proves that the actual integral
idempotent $e$ is self-adjoint:
\[
                         \langle e\xi,\eta\rangle
                                  =\langle\xi,e\eta\rangle .
                                                               \tag{D.70p}
\]
Hence (D.70m) is an orthogonal direct sum.
Each restricted form is unimodular: in bases
of the two integral direct summands, the whole
Gram matrix is block diagonal with unit
determinant, so each block determinant is a
unit in $\Z_2$.
In particular the rank-$2q-2$ form on $W_2$
is nonzero and unimodular.
Choose $\theta,\psi\in W_2$ with
$\langle\theta,\psi\rangle$ a unit.
Their reductions have nonzero cup product
in $H^2(P,\F_2)$, while (D.70n) makes both
characters fixed by the actual $b$.
The degree-three restriction obstruction
(D.70a)--(D.70g) applies to exactly this $P,b$,
and says that their reduced cup product is zero.
This contradiction excludes the chosen
coordinate $\lambda_q=-1$.

There is no negative odd coordinate.
Together with $\lambda^2=1$ and $\lambda_2=1$,
this proves $\lambda=1$ in the full product
$\widehat\Z^\times=\prod_p\Z_p^\times$.
The possible finite two-subgroups are exactly
$1,C_2,C_4,C_2\times C_2,D_8$ by D.50.
D.43 gives the global sign for their involutions;
the present result gives multiplier one for
their order-four elements. These statements
cover every element of each of the listed
groups and prove the final assertion.
\end{proof}

\emph{Scope.}
This result closes the quartic multiplier branch
left in D.50 and D.68. It uses the exact
rotation-cover conclusion of D.69 and the actual
closed filling (D.70i).
It supplies no conjugacy to a geometric quartic
representative and no recognition of either full
fixed subgroup in D.69.
The remaining scalar order-three branch is
unaffected by this degree-three calculation.
